\section{Full KDE4 Production Run}
\label{sec:kde4run}

The benchmarks in Chapters~\ref{sec:batching} and~\ref{sec:placeholders} timed the model alone on samples of up to 1,000
strings. This chapter puts all the changes into the \textbf{production pipeline} (Phases~1--3, Translation Memory, S3,
telemetry) and translates the \textbf{whole KDE4 dataset}.

\subsection{Setup}
All 10,054 usable KDE4 strings were packed into one \code{.resx} file and translated into all 12 languages:
\textbf{120,648 translations}. The pipeline used the adopted batching configuration (Table~\ref{tab:adopted}), beam~2 and
placeholder protection, on the same m5.large. The Translation Memory was emptied first, so every string went through the
model (run \code{ec2-20261016t214343z}).

\subsection{Results}
The run succeeded: 12 of 12 language files, \textbf{0 failures, 0 missing strings}.

\begin{table}[H]
\centering
\caption[Full KDE4 run per language: speed]{Full KDE4 run per language: speed and placeholder survival, 10,054 strings per language.}
\label{tab:kde4run}
\begin{tabular}{@{}l S[table-format=2.2] S[table-format=5.0] S[table-format=3.2]@{}}
\toprule
\textbf{Language} & {\textbf{Time (min)}} & {\textbf{Translations/hour}} & {\textbf{Placeholders kept (\%)}} \\
\midrule
Arabic (ar) & 18.28 & 32991 & 99.72 \\
Bengali (bn) & 18.38 & 32828 & 100.00 \\
German (de) & 17.83 & 33841 & 99.44 \\
Spanish (es) & 17.62 & 34232 & 100.00 \\
French (fr) & 19.00 & 31742 & 99.72 \\
Hindi (hi) & 16.12 & 37429 & 100.00 \\
Japanese (ja) & 16.18 & 37282 & 100.00 \\
Marathi (mr) & 17.02 & 35434 & 100.00 \\
Punjabi (pa) & 19.30 & 31258 & 100.00 \\
Tamil (ta) & 18.77 & 32144 & 100.00 \\
Telugu (te) & 19.15 & 31502 & 100.00 \\
Chinese (zh) & 16.73 & 36059 & 99.72 \\
\midrule
\textbf{All} & \bfseries 214.4 & \bfseries 33760 & \bfseries 99.88 \\
\bottomrule
\end{tabular}
\end{table}

\begin{table}[H]
\centering
\caption[Translation quality of the full KDE4 run]{Translation quality of the full KDE4 run. chrF, BLEU and spBLEU use
all 10,054 strings per language; COMET uses the same fixed sample of 1,000 strings in every language.}
\label{tab:kde4quality}
\begin{tabular}{@{}l S[table-format=2.2] S[table-format=2.2] S[table-format=2.2] S[table-format=1.4]@{}}
\toprule
\textbf{Language} & {\textbf{chrF}} & {\textbf{BLEU}} & {\textbf{spBLEU}} & {\textbf{COMET}} \\
\midrule
Arabic (ar) & 55.42 & 32.45 & 39.44 & 0.8296 \\
Bengali (bn) & 64.38 & 37.38 & 52.52 & 0.8866 \\
German (de) & 64.22 & 39.92 & 45.88 & 0.8639 \\
Spanish (es) & 71.50 & 50.65 & 53.15 & 0.8704 \\
French (fr) & 68.89 & 44.18 & 46.65 & 0.8495 \\
Hindi (hi) & 64.75 & 43.24 & 49.75 & 0.8526 \\
Japanese (ja) & 50.92 & 32.35 & 39.07 & 0.8643 \\
Marathi (mr) & 58.45 & 35.76 & 44.38 & 0.8249 \\
Punjabi (pa) & 61.53 & 38.56 & 50.42 & 0.8689 \\
Tamil (ta) & 64.10 & 37.75 & 49.53 & 0.8840 \\
Telugu (te) & 65.28 & 40.85 & 53.82 & 0.8932 \\
Chinese (zh) & 39.84 & 42.48 & 42.55 & 0.8530 \\
\midrule
\textbf{All} & \bfseries 63.29 & {\bfseries 39.63$^{\ast}$} & \bfseries 47.90 & \bfseries 0.8617 \\
\bottomrule
\end{tabular}

\smallskip
{\footnotesize $^{\ast}$Mean of the 12 languages: BLEU uses different word splitting for Chinese and Japanese, so a single
corpus score across languages is not meaningful. spBLEU uses the same pieces for every language.}
\end{table}

\begin{figure}[H]
\centering
\begin{tikzpicture}
\begin{axis}[reportbars, width=\linewidth, height=5.2cm, bar width=12pt,
  symbolic x coords={ar,bn,de,es,fr,hi,ja,mr,pa,ta,te,zh}, xtick=data,
  ylabel={COMET}, ymin=0.78, ymax=0.91, enlarge x limits=0.05,
  yticklabel style={/pgf/number format/fixed, /pgf/number format/precision=2},
  nodes near coords, nodes near coords style={font=\tiny, rotate=90, anchor=west, /pgf/number format/fixed, /pgf/number format/fixed zerofill, /pgf/number format/precision=3}]
\addplot[fill=cnew!80, draw=cnew] coordinates {(ar,0.8296) (bn,0.8866) (de,0.8639) (es,0.8704) (fr,0.8495) (hi,0.8526) (ja,0.8643) (mr,0.8249) (pa,0.8689) (ta,0.8840) (te,0.8932) (zh,0.8530)};
\end{axis}
\end{tikzpicture}
\caption[COMET per language in the full KDE4 run]{COMET per language in the full KDE4 run (1,000-string sample per
language). Every language lies between 0.82 and 0.90.}
\label{fig:kde4chrf}
\end{figure}

\paragraph{Reading the metrics together.} The four metrics agree on most languages, with one exception. chrF places
Chinese far below every other language (39.84), but BLEU (42.48), spBLEU (42.55) and COMET (0.853) put it in the middle.
chrF compares sequences of up to six characters, which is harsh for Chinese, where a word is often one or two
characters; COMET, which judges meaning with a neural model, is the fairer comparison across scripts.

\paragraph{Where the time goes.} Translation took \textbf{3~h~34~min} (3~h~38~min for the whole run, including model
loading and Phases~1 and~3). Model inference accounted for 99.6\% of Phase~2; Translation Memory reads and writes took
about 44~s in total, and XML and S3 a few seconds. During translation the machine's CPU was 92--100\% busy, and memory use
stayed at 1.6--1.9~GB per file.

\subsection{Comparison with the Starting Configuration}
The starting configuration (beam~5, one batch at a time) ran at 13,631 translations per hour in the KDE4 benchmark
(Table~\ref{tab:kde4beam}). At that rate this dataset would take an estimated \textbf{8.9 hours}; the optimized
production pipeline needed \textbf{3.6 hours}, about \textbf{2.5\x{} faster}. The benchmarks predicted 2.8\x{}; the
difference is the work the benchmark does not time: tokenization, placeholder handling, XML, database saves, and the end
of each 1,000-string chunk, where only one batch is left for the two cores.

\keyresult{\textbf{RQ3.} The optimized pipeline translated the whole twelve-language KDE4 corpus (120,648 translations)
in 3.6 hours with no failures, at 33,760 translations per hour, with chrF 63.29, spBLEU 47.90, COMET 0.862 and
99.88\% of placeholders kept: about 2.5\x{} faster than the semester's starting configuration.}
